\begin{table}[!htbp]
\centering
\caption{The full event window and a common initial seven-day window}\label{tab:timing}
\small
\setlength{\tabcolsep}{5pt}
\begin{tabular}{llrrrrr}
\toprule
Type & Period & Events & Gap & Full & Initial seven & Remainder \\
\midrule
Addition & 2007--2011 & 33 & 14 & 8.52 & 1.19 & 7.34 \\
Addition & 2012--2016 & 54 & 16 & 0.25 & 1.26 & -1.01 \\
Addition & 2017--2021 & 82 & 15 & 2.85 & 1.23 & 1.62 \\
Addition & 2022--2026 & 87 & 6 & 1.50 & 1.80 & -0.31 \\
Deletion & 2007--2011 & 33 & 15 & -6.16 & -3.54 & -2.62 \\
Deletion & 2012--2016 & 46 & 16 & -2.47 & -0.86 & -1.61 \\
Deletion & 2017--2021 & 61 & 15 & -0.43 & -1.82 & 1.39 \\
Deletion & 2022--2026 & 75 & 6 & -1.53 & -0.76 & -0.77 \\
\bottomrule
\end{tabular}
\begin{minipage}{\linewidth}\footnotesize\vspace{5pt}
Returns are percentages. Gap is the median AD-to-ED distance in trading days across distinct reviews represented in the sample. Full includes AD-1 through ED+1; Initial seven includes AD-1 through AD+5; Remainder includes AD+6 through ED+1. Every row uses identical events and the two components sum to Full. The remainder has varying duration. This timing diagnostic is exploratory; uncertainty for period contrasts is reported in Appendix~\ref{app:inference}.
\end{minipage}
\end{table}
